\clearpage\makeatletter\@mainmatterfalse\makeatother

% BEGIN bibliography.tex
\renewcommand{\bibname}{第一卷原著参考文献}
\begin{thebibliography}{99}
\addcontentsline{toc}{chapter}{第一卷原著参考文献}
\markboth{第一卷原著参考文献}{第一卷原著参考文献}
\small
\bibitem[1]{b1}
H. Cartan 与 C. Chevalley：巴黎高等师范学校讨论班，第 8 年（1955—1956），代数几何。
\textit{Séminaire de l'École Normale Supérieure}, 8\textsuperscript{e} année (1955–56),
\textit{Géométrie algébrique}.

\bibitem[2]{b2}
H. Cartan 与 S. Eilenberg：\textit{Homological Algebra}（《同调代数》），
Princeton Math. Series，Princeton University Press，1956。

\bibitem[3]{b3}
W. L. Chow 与 J. Igusa：Cohomology theory of varieties over rings（环上簇的上同调理论），
\textit{Proc. Nat. Acad. Sci. U.S.A.}，第 XLIV 卷（1958），1244–1248 页。

\bibitem[4]{b4}
R. Godement：\textit{Théorie des faisceaux}（《层论》），
Actual. Scient. et Ind.，第 1252 号，Paris：Hermann，1958。

\bibitem[5]{b5}
H. Grauert：Ein Theorem der analytischen Garbentheorie und die Modulräume komplexer Strukturen
（解析层论的一个定理与复结构的模空间），
\textit{Publ. Math. Inst. Hautes Études Scient.}，第 5 号，1960。

\bibitem[6]{b6}
A. Grothendieck：Sur quelques points d'algèbre homologique（同调代数的若干问题），
\textit{Tôhoku Math. Journ.}，第 IX 卷（1957），119–221 页。

\bibitem[7]{b7}
A. Grothendieck：Cohomology theory of abstract algebraic varieties（抽象代数簇的上同调理论），
\textit{Proc. Intern. Congress of Math.}，103–118 页，Edinburgh（1958）。

\bibitem[8]{b8}
A. Grothendieck：Géométrie formelle et géométrie algébrique（形式几何与代数几何），
\textit{Séminaire Bourbaki}，第 11 年（1958—1959），第 182 讲。

\bibitem[9]{b9}
M. Nagata：A general theory of algebraic geometry over Dedekind domains
（Dedekind 整环上代数几何的一般理论），\textit{Amer. Math. Journ.}：
I，第 LXXVIII 卷，78–116 页（1956）；II，第 LXXX 卷，382–420 页（1958）。

\bibitem[10]{b10}
D. G. Northcott：\textit{Ideal theory}（《理想论》），Cambridge University Press，1953。

\bibitem[11]{b11}
P. Samuel：\textit{Commutative algebra}（《交换代数》；D. Herzig 记录），Cornell University，1953。

\bibitem[12]{b12}
P. Samuel：\textit{Algèbre locale}（《局部代数》），Mém. Sci. Math.，第 123 号，Paris，1953。

\bibitem[13]{b13}
P. Samuel 与 O. Zariski：\textit{Commutative algebra}（《交换代数》），两卷，
New York：Van Nostrand，1958–60。

\bibitem[14]{b14}
J.-P. Serre：Faisceaux algébriques cohérents（凝聚代数层），
\textit{Ann. of Math.}，第 LXI 卷（1955），197–278 页。

\bibitem[15]{b15}
J.-P. Serre：Sur la cohomologie des variétés algébriques（论代数簇的上同调），
\textit{Journ. de Math.} (9)，第 XXXVI 卷（1957），1–16 页。

\bibitem[16]{b16}
J.-P. Serre：Géométrie algébrique et géométrie analytique（代数几何与解析几何），
\textit{Ann. Inst. Fourier}，第 VI 卷（1955–56），1–42 页。

\bibitem[17]{b17}
J.-P. Serre：Sur la dimension homologique des anneaux et des modules noethériens
（论诺特环与诺特模的同调维数），
\textit{Proc. Intern. Symp. on Alg. Number Theory}，176–189 页，Tokyo—Nikko，1955。

\bibitem[18]{b18}
A. Weil：\textit{Foundations of algebraic geometry}（《代数几何基础》），
Amer. Math. Soc. Coll. Publ.，第 29 号，1946。

\bibitem[19]{b19}
A. Weil：Numbers of solutions of equations in finite fields（有限域中方程的解的数目），
\textit{Bull. Amer. Math. Soc.}，第 LV 卷（1949），497–508 页。

\bibitem[20]{b20}
O. Zariski：\textit{Theory and applications of holomorphic functions on algebraic varieties over arbitrary ground fields}
（任意基域上的代数簇的全纯函数理论及应用），Mem. Amer. Math. Soc.，第 5 号（1951）。

\bibitem[21]{b21}
O. Zariski：A new proof of Hilbert's Nullstellensatz（Hilbert 零点定理的一个新证明），
\textit{Bull. Amer. Math. Soc.}，第 LIII 卷（1947），362–368 页。

\bibitem[22]{b22}
E. Kähler：Geometria Arithmetica（算术几何），\textit{Ann. di Mat.} (4)，第 XLV 卷（1958），1–368 页。
\end{thebibliography}


% END bibliography.tex


% BEGIN notation_index.tex
\chapter*{第一卷原著记号索引}
\addcontentsline{toc}{chapter}{第一卷原著记号索引}
\markboth{第一卷原著记号索引}{第一卷原著记号索引}
\src{217—218}
\noindent\textbf{译注。}以下按原著次序译出全部索引条目。右栏是原索引所列的章、条目号，不是本译稿页码；原索引中个别沿用旧编号的引文仍予保留，不静默改写。查阅时须与正文条目核对。
\begingroup\small
\renewcommand{\arraystretch}{1.15}
\begin{longtable}{@{}p{0.78\linewidth}p{0.17\linewidth}@{}}
\toprule 记号及说明&原索引引文\\\midrule\endfirsthead
\toprule 记号及说明（续）&原索引引文\\\midrule\endhead
\bottomrule\endfoot
$M_{[\varphi]}$（$M$ 为 $B$-模，$\varphi:A\to B$ 为环同态）&0，1.0.2\\
$BJ,JB$（$J$ 为环 $A$ 的理想，已给定 $A\to B$ 的同态）&0，1.0.3\\
$r(\mathfrak a)$（$\mathfrak a$ 为理想）&0，1.1.1\\
$\mathfrak R(A)$（$A$ 为环）&0，1.1.2\\
$S_f$（$f$ 为交换环 $A$ 的元素）&0，1.2.1\\
$S^{-1}A$，$S^{-1}M$，$m/s$，$i_A^S$，$i_M^S$，$i^S$（$A$ 为环，$M$ 为 $A$-模，$S$ 为 $A$ 的乘法子集）&0，1.2.2\\
$A_f,M_f,A_{\mathfrak p},M_{\mathfrak p}$（$M$ 为 $A$-模，$f\in A$，$\mathfrak p$ 为 $A$ 的素理想）&0，1.2.3\\
$S^{-1}u$（$u$ 为 $A$-模同态，$S$ 为 $A$ 的乘法子集）&0，1.3.1\\
$\rho_A^{T,S},\rho_M^{T,S},\rho^{T,S}$（$M$ 为 $A$-模，$S,T$ 为 $A$ 的乘法子集）&0，1.4.1\\
$\Supp(M)$（$M$ 为 $A$-模）&0，1.7.1\\
$\F|U,u|U$（$\F$ 为 $X$ 上的层，$u$ 为 $X$ 上的层态射，$U$ 为 $X$ 的开集）&0，3.1.5\\
$\F_x,s_x,\Gamma(U,\F),u(s),\Supp(\F)$（$\F$ 为 $X$ 上的集合层，$x\in X$，$U$ 为开集，$s\in\Gamma(U,\F)$，$u$ 为层同态）&0，3.1.6\\
$\psi_*(\F)$（$\F$ 为 $X$ 上的层，$\psi:X\to Y$ 连续）&0，3.4.1\\
$\psi_*(u)$（$u$ 为 $X$ 上的层同态）&0，3.4.2\\
$\psi_x$&0，3.4.4\\
$\G\to\F$（$\F$ 为 $X$ 上的层，$\G$ 为 $Y$ 上的层）&0，3.5.1\\
$u^\sharp,v^\flat,\rho_{\G}$&0，3.5.3\\
$\psi^*(\G),\psi^*(v),\sigma_{\F}$&0，3.5.5\\
$\OO_X,\OO_{X,x},\OO_x,1,e$（$X$ 为赋环空间）&0，4.1.1\\
$\dual\F,\bigwedge^p\F$（$\F$ 为 $\OO_X$-模层）&0，4.1.4\\
$\J\F$（$\J$ 为 $\OO_X$ 的理想层，$\F$ 为 $\OO_X$-模层）&0，4.1.5\\
$\Psi_*(\F),\Psi_*(u)$（$\F$ 为 $\OO_X$-模层，$u$ 为 $\OO_X$-模层同态）&0，4.2.1\\
$\Psi_*(\C)$（$\C$ 为 $\OO_X$-代数层）&0，4.2.4\\
$\Psi^*(\G),\Psi^*(v)$（$\G$ 为 $\OO_Y$-模层，$v$ 为 $\OO_Y$-模层同态）&0，4.3.1\\
$\Psi^*(\C)$（$\C$ 为 $\OO_Y$-代数层）&0，4.3.4\\
$\Psi^*(\J)\OO_X,\J\OO_X$（$\J$ 为 $\OO_Y$ 的理想）&0，4.3.5\\
$\G\to\F$（$\F$ 为 $\OO_X$-模层，$\G$ 为 $\OO_Y$-模层）&0，4.4.1\\
$u^\sharp_\Psi,u^\sharp,v^\flat_\Psi,v^\flat,\rho_{\G},\sigma_{\F}$&0，4.4.3\\
$u_1\otimes u_2$（$u_1,u_2$ 为从 $\OO_Y$-模层到 $\OO_X$-模层的同态）&0，4.4.4\\
$\LL^{-1}$（$\LL$ 为可逆 $\OO_X$-模层）&0，5.4.3\\
$\LL^{\otimes n}$（$\LL$ 为可逆 $\OO_X$-模层）&0，5.4.4\\
$\Gamma_*(\LL),\Gamma_*(\LL,\F)$（$\LL$ 可逆，$\F$ 为 $\OO_X$-模层）&0，5.4.6\\
$\OO_X^*$&0，5.4.7\\
$\mathfrak m_x,k(x),f(x)$&0，5.5.1\\
$X_f$&0，5.5.2\\
$\Im(M'\otimes_A N')$（$M',N'$ 分别为 $M,N$ 的子 $A$-模）&0，6.0\\
$\widehat A,\widehat M$&0，7.2.3及7.3.1\\
$A\{T_1,\ldots,T_r\}$&0，7.5.1\\
$A\{S^{-1}\}$&0，7.6.1\\
$\mathfrak a\{S^{-1}\}$（$\mathfrak a$ 为 $A$ 的开理想）&0，7.6.9\\
$A_{\{f\}},\mathfrak a_{\{f\}}$&0，7.6.14\\
$A_{\{S\}}$&0，7.6.14\\
$\widehat{M\otimes_A N},M\ctensor_A N$&0，7.7.1\\
$u\ctensor v$&0，7.7.3\\
\midrule
$\Spec(A),\jj_x,\mathfrak m_x,k(x),f(x),M_x$，$r(E),V(E),V(f),D(f)$（$A$ 为环，$M$ 为 $A$-模，$f\in A$，$E\subset A$，$x\in\Spec(A)$）&I，1.1.1\\
$\jj(Y)$（$Y\subset\Spec(A)$）&I，1.1.3\\
${}^{a}\varphi$（$\varphi$ 为环同态）&I，1.2.1\\
$S_f'$（$f$ 为环的元素）&I，1.3.1\\
$\rho_{g,f}$（$f,g$ 为环的元素）&I，1.3.3\\
$\wt A,\wt M,\theta_f$（$A$ 为环，$f\in A$，$M$ 为 $A$-模）&I，1.3.4\\
$\wt u$（$u$ 为 $A$-模同态）&I，1.3.5\\
$\wt\varphi$（$\varphi$ 为环同态）&I，1.6.1\\
$A(X)$（$X$ 为仿射概形）&I，1.7.1\\
$\OO_{X/Y}$（$X$ 为预概形）&I，2.1.6\\
$\Hom(X,Y)$（$X,Y$ 为预概形）&I，2.2.1\\
$\Hom_S(X,Y),1_X$（$X,Y$ 为 $S$-预概形）&I，2.5.2\\
$\Gamma(X/S)$（$S$ 为预概形，$X$ 为 $S$-预概形）&I，2.5.5\\
$X\amalg Y$（$X,Y$ 为预概形）&I，3.1.1\\
$X\times_S Y,X\times Y,(g,h)_S,u\times_Sv,u\times v$（$X,Y$ 为 $S$-预概形，$g,h,u,v$ 为 $S$-态射）&I，3.2.1\\
$X\times_A Y,X\otimes_A B,(g,h)_A,u\times_Av$（$A$ 为环，$X,Y$ 为 $A$-预概形，$B$ 为 $A$-代数，$g,h,u,v$ 为 $A$-态射）&I，3.2.1\\
$X_{(S')}$（$X,S'$ 为 $S$-预概形）&I，3.3.6\\
$f_{(S')}$（$S'$ 为 $S$-预概形，$f$ 为 $S$-态射）&I，3.3.7\\
$\Gamma_f$（$f$ 为 $S$-态射）&I，3.3.14\\
$X(T)$（$X,T$ 为预概形）&I，3.4.1\\
$P\times_R Q$（$P,Q$ 为 $R$ 上的集合）&I，3.4.2\\
$X(T)_S$（$X,T$ 为 $S$-预概形）&I，3.4.3\\
$X(B),X(B)_A$（$X$ 为 $A$-预概形，$B$ 为 $A$-代数）&I，3.4.4\\
$X\otimes_Y B,X\otimes_{\OO_y}B$（$y\in Y$，$B$ 为 $\OO_y$-代数）&I，3.6.3\\
$Z\le Y$（$Y,Z$ 为一个预概形的子预概形）&I，4.1.10\\
$f^{-1}(Y)$（$f$ 为态射，$Y$ 为子预概形）&I，4.4.1\\
$\mathscr N_X$（$X$ 为预概形）&I，5.1.1\\
$X_{\mathrm{red}}$（$X$ 为预概形）&I，5.1.3\\
$f_{\mathrm{red}}$（$f$ 为态射）&I，5.1.5\\
$\Delta_{X/S},\Delta_X,\Delta$（$X$ 为 $S$-预概形）&I，5.3.1\\
$\operatorname{rg}_K(X)$（$K$ 为域，$X$ 为有限 $K$-概形）&I，6.4.5\\
$n(X)$（$X$ 为域上的有限概形）&I，6.4.8\\
$\Gamma_{\mathrm{rat}}(X/Y)$&I，7.1.2\\
$R(X)$（$X$ 为预概形）&I，7.1.3\\
$\mathscr R(X)$（$X$ 为预概形）&I，7.3.1\\
$L(A)$（$A$ 为整环）&I，8.1.2\\
$\delta(f)$（$f$ 为有理映射）&I，8.2.1\\
$\F\otimes_{\OO_S}\G,\F\otimes_S\G$（$\F,\G$ 为 $S$-预概形上的模层）&I，9.1.2\\
$\overline\G$（$\G$ 为 $\OO_X$-模层）&I，9.4.1\\
$\overline Y$（$Y$ 为子预概形）&I，9.5.11\\
$\Spf(A),\OO_{\frX}$（$A$ 为可容许环，$\frX=\Spf(A)$）&I，10.1.2\\
$\frD(f)$（$f$ 为可容许环的元素）&I，10.1.4\\
${}^{a}\varphi,\wt\varphi$（$\varphi$ 为可容许环的连续同态）&I，10.2.1\\
$J^\Delta$（$J$ 为定义理想）&I，10.3.1\\
$\frX\times_{\frS}\frY$（$\frX,\frY$ 为 $\frS$-形式预概形）&I，10.7.3\\
$\F_{/X'},\widehat\F,u_{/X'},\widehat u$（$\F$ 为 $\OO_X$-模层，$u$ 为模层同态，$X'$ 为 $X$ 的闭子集）&I，10.8.4\\
$X_{/X'},\widehat X$（$X$ 为预概形，$X'$ 为闭子集）&I，10.8.5\\
$\widehat f$（$f$ 为预概形态射）&I，10.9.1\\
$M^\Delta,u^\Delta$（$M$ 为进制环 $A$ 上的模，$u$ 为 $A$-模的连续同态）&I，10.10.1\\
$\Delta_{\frX/\frS},\Delta_{\frX}$（$\frX$ 为 $\frS$-形式预概形）&I，10.15.1\\
\end{longtable}
\endgroup


% END notation_index.tex

\par

% BEGIN terminology_index.tex
\chapter*{第一卷原著术语索引}
\addcontentsline{toc}{chapter}{第一卷原著术语索引}
\markboth{第一卷原著术语索引}{第一卷原著术语索引}
\src{219—224}
\noindent\textbf{译注。}本索引依原著法文词条的次序逐条翻译；右栏保留原著所列的章号及条目号，不是中文页码。原索引有若干旧编号或误引，例如“闭点”所引的（0，2.2.6）在本卷正文中并不存在；这里不静默改动原引文。查阅时请同时参照本书目录和正文。
\begingroup\small
\renewcommand{\arraystretch}{1.16}
\begin{longtable}{@{}p{0.76\linewidth}p{0.19\linewidth}@{}}
\toprule 术语&原索引引文\\\midrule\endfirsthead
\toprule 术语（续）&原索引引文\\\midrule\endhead
\bottomrule\endfoot
子预概形的闭包&I，9.5.11\\
$\OO_X$-代数层&0，4.1.3\\
进制环；$J$-进制环&0，7.1.9\\
可容许环&0，7.1.2\\
完备分式环&0，7.6.5\\
分式环&0，1.2.2\\
有理函数环&I，2.1.6；7.1.3\\
仿射概形的环&I，1.7.1\\
整环&0，1.0.6\\
线性拓扑环&0，7.1.1\\
局部环&0，1.0.7\\
支配局部环&I，8.1.1\\
$X$ 沿 $Y$ 的局部环；$Y$ 在 $X$ 中的局部环&I，2.1.6\\
预进制环；$J$-预进制环&0，7.1.9\\
预可容许环&0，7.1.2\\
约化环&0，1.1.1\\
正则环&0，4.1.3\\
相关的局部环&I，8.1.4\\
$\OO_X$-模层的零化理想&0，5.3.7\\
与环同态相伴的环谱映射&I，1.2.1\\
有理映射；$S$-有理映射&I，7.1.2\\
在一点有定义的有理映射&I，7.2.1\\
在开集上诱导的有理映射&I，7.1.2\\
在 $\Spec(\OO_x)$ 上诱导的有理映射&I，7.2.8\\
凝聚的 $\OO_X$-模层&0，5.3.1\\
凝聚的 $\OO_X$-代数层&0，5.3.6\\
$\OO_X$-模层及其同态沿闭子集的完备化&I，10.8.4\\
预概形沿闭子集的完备化&I，10.8.5\\
不可约分支&0，2.1.5\\
$\psi$-态射与 $\psi'$-态射的复合&0，3.5.2\\
$\Psi$-态射与 $\Psi'$-态射的复合&0，4.4.2\\
粘合条件&0，3.3.1；4.1.6\\
代数预概形的基域&I，6.4.1\\
几何点的值域&I，3.4.5\\
$X\times_S X$ 的对角线&I，5.3.9\\
双同态&0，1.0.2\\
有理映射的定义域&I，7.2.1\\
$\OO_X$-模层的对偶&0，4.1.4\\
拓扑幂零元素&0，7.1.1\\
由一族截面生成的 $\OO_X$-模层&0，5.1.2\\
截面为零的点集&0，5.5.1\\
整代数；有限整代数&0，1.0.5\\
赋环空间&0，4.1.1\\
局部环赋环空间&0，5.5.1\\
在开集上诱导的赋环空间&0，4.1.2\\
正规的、约化的、正则的赋环空间&0，4.1.3\\
由粘合得到的赋环空间&0，4.1.6\\
柯尔莫哥洛夫空间&0，2.1.2\\
不可约空间&0，2.1.1\\
诺特空间&0，2.2.1\\
拟紧空间&0，2.2.4\\
赋环空间的底空间&0，4.1.1\\
拓扑赋环空间&0，4.1.1\\
赋环空间上的代数层&0，4.1.3\\
与预层相伴的层&0，3.5.7\\
$\Spec(A)$ 上与一个 $A$-模相伴的层&I，1.3.4\\
取值于一个范畴的层&0，3.1.1\\
凝聚环层&0，5.3.5\\
在一点正规的、正规的、在一点约化的、约化的、在一点正则的、正则的环层&0，4.1.3\\
分次环层&0，4.1.3\\
有理函数层&I，7.3.1\\
挠层&I，7.4.1\\
理想层&0，4.1.3\\
定义理想层&I，10.3.3；10.5.1\\
诱导层&0，3.7.1\\
局部常值层&0，3.6.1\\
由粘合得到的层&0，3.3.1\\
伪离散层&0，3.8.1\\
常值层&0，3.6.1\\
赋环空间的结构层&0，4.1.1\\
仿射概形的结构层&I，1.3.4\\
开集基上的层&0，3.2.2\\
层的茎&0，3.1.6\\
预概形态射的纤维&I，3.6.3\\
有限代数&0，1.0.5\\
有理函数&I，7.1.2\\
一点的一般化&0，2.1.2\\
分次 $\OO_X$-模层&0，4.1.3\\
态射的图&I，5.3.11\\
拓扑环层的连续同态&0，3.1.4\\
由截面定义的同态&0，5.1.1\\
局部环的局部同态&0，1.0.7\\
模的 $\varphi$-同态&0，1.0.2\\
可容许环的定义理想&0，7.1.2\\
$\OO_X$-理想&0，4.1.3\\
形式预概形 $\frX$ 的 $\OO_{\frX}$-定义理想&I，10.3.3；10.5.1\\
素理想&0，1.0.6\\
$\OO_X$-模层的直像&0，4.2.1\\
预层的直像&0，3.4.1\\
$S$-截面的像&I，5.3.11\\
预概形在态射下的闭像&I，9.5.3\\
$\OO_Y$-模层的逆像&0，4.3.1\\
$S$-态射的逆像&I，3.3.7\\
预层的逆像&0，3.5.3\\
$S$-预概形的逆像&I，3.3.6\\
子预概形的逆像&I，4.4.1\\
预概形的浸入、闭浸入、开浸入&I，4.2.1\\
形式预概形的闭浸入&I，10.14.2\\
预概形的局部浸入&I，4.5.1\\
开集上诱导的赋环空间的典范嵌入&0，4.1.2\\
子预概形的典范嵌入&I，4.1.7\\
形式子预概形的典范嵌入&I，10.14.2\\
几何单射&I，3.5.4\\
可逆 $\OO_X$-模层的逆&0，5.4.3\\
可逆 $\OO_X$-模层&0，5.4.1\\
与浸入相伴的同构&I，4.2.1\\
预概形的局部同构&I，4.5.1\\
$\OO_{X_n}$-模层的射影极限&I，10.6.6\\
局部自由 $\OO_X$-模层&0，5.4.1\\
取值于局部环的点的落点&I，3.4.5\\
分式模&0，1.2.2\\
忠实平坦模&0，6.4.1\\
平坦模&0，6.1.1\\
$A$-平坦模&0，6.2\\
拟有限模&0，7.4.1\\
$\OO_X$-模层&0，4.1.3\\
赋环空间、拓扑赋环空间的态射&0，4.1.1\\
忠实平坦态射&0，6.7.8\\
平坦态射&0，6.7.1\\
开集基上所定义预层的态射&0，3.2.3\\
预概形态射&I，2.2.1\\
$S$-预概形、$A$-预概形的态射&I，2.5.2\\
双有理态射&I，6.5.6\\
有限型态射&I，6.3.1\\
对角态射&I，5.3.1\\
支配态射&I，2.2.6\\
闭态射&I，2.2.6\\
一个态射的图态射&I，3.3.14\\
局部有限型态射&I，6.6.2\\
不大于另一态射的态射&I，4.1.8\\
开态射&I，2.2.6\\
拟紧态射&I，6.6.1\\
根式态射&I，3.5.4\\
约化态射&I，5.1.5\\
分离态射&I，5.4.1\\
$S$-预概形的结构态射&I，2.5.1\\
满态射&I，2.2.6\\
普遍单射态射&I，3.5.4\\
等价的态射&I，7.1.1\\
预概形的 $A$-态射、$S$-态射&I，2.5.2\\
形式预概形态射&I，10.4.5\\
形式预概形的进制态射&I，10.12.1\\
形式预概形的对角态射&I，10.15.1\\
形式预概形的有限型态射&I，10.13.3\\
形式预概形的分离态射&I，10.15.1\\
$\frS$-形式预概形的结构态射&I，10.4.7\\
形式预概形的 $A$-态射、$\frS$-态射&I，10.4.7\\
预层的 $\psi$-态射&0，3.5.1\\
从 $\OO_Y$-模层到 $\OO_X$-模层的 $\Psi$-态射&0，4.4.1\\
环的幂零根&0，1.1.1\\
$\OO_X$-代数层的幂零根&I，5.1.1\\
预概形的几何点数&I，6.4.8\\
仿射开集&I，2.1.1\\
仿射形式开集、进制仿射形式开集、诺特仿射形式开集&I，10.4.1\\
环的乘法子集&0，1.2.1\\
饱和乘法子集&0，1.4.3\\
$f$-平坦的 $\OO_X$-模层&0，6.7.1\\
预概形的环值点&I，3.4.4\\
预概形的预概形值点&I，3.4.1\\
$A$-预概形的 $A$-代数值点&I，3.4.4\\
$S$-预概形在 $s\in S$ 上方的点&I，2.5.1\\
$S$-预概形的 $S$-预概形值点&I，3.4.3\\
闭点&0，2.2.6\\
泛点&0，2.1.2\\
预概形的几何点&I，3.4.5\\
$s$ 上方的几何点&I，3.4.5\\
落在 $x$ 的几何点&I，3.4.5\\
$K$ 上的有理点&I，3.4.5\\
常值预层&0，3.6.1\\
开集基上的预层&0，3.2.1\\
预概形&I，2.1.2\\
阿廷预概形&I，6.2.1\\
连通预概形&I，2.1.7\\
基预概形&I，2.5.1\\
模 $J$ 约化所得的预概形&I，3.7.1\\
在开集上诱导的预概形&I，2.1.8\\
整预概形&I，2.1.7\\
不可约预概形&I，2.1.7\\
局部整预概形&I，2.1.7\\
局部诺特预概形&I，6.1.1\\
与预概形相伴的约化预概形&I，5.1.3\\
$A$-预概形；$A$ 上的预概形&I，2.5.1\\
代数 $K$-预概形&I，6.4.5\\
$S$-预概形；$S$ 上的预概形&I，2.5.1\\
支配的 $S$-预概形&I，2.5.1\\
$S$ 上有限型预概形；有限型 $S$-预概形&I，6.3.1\\
由基预概形扩张得到的预概形&I，3.3.6\\
在 $S$ 上分离的预概形&I，5.4.1\\
形式预概形&I，10.4.2\\
进制形式预概形&I，10.4.2\\
局部诺特形式预概形&I，10.4.2\\
诺特形式预概形&I，10.4.2\\
$A$ 上的形式预概形；$A$-形式预概形&I，10.4.7\\
$\frS$ 上的形式预概形；$\frS$-形式预概形&I，10.4.7\\
$\frS$ 上有限型形式预概形；有限型 $\frS$-形式预概形&I，10.13.3\\
在 $S$ 上分离的形式预概形&I，10.15.1\\
具有有限表示的模&0，1.0.5\\
具有有限表示的 $\OO_X$-模层&0，5.2.5\\
诺特归纳原理&0，2.2.2\\
$S$-预概形的积&I，3.2.1\\
$\frS$-形式预概形的积&I，10.7.1\\
集合的纤维积&I，3.4.2\\
不同预概形上层的张量积&I，9.1.2\\
代数的完备张量积&0，7.7.5\\
模的完备张量积&0，7.7.2\\
积的典范投影&I，3.2.1\\
子模层的典范延拓&I，9.4.1\\
态射到完备化的延拓&I，10.9.1\\
$\OO_X$-模层的第 $p$ 外幂&0，4.1.4\\
拟凝聚的 $\OO_X$-模层&0，5.1.3\\
拟凝聚的 $\OO_X$-代数层&0，5.1.3\\
理想的根&0，1.1.1\\
环的根&0，1.1.2\\
局部自由 $\OO_X$-模层的秩&0，5.4.1\\
无挠 $\OO_X$-模层的秩&I，7.4.2\\
有限 $K$-概形的秩&I，6.4.5\\
有限 $K$-概形的可分秩&I，6.4.8\\
赋环空间在开集上的限制&0，4.1.2\\
态射在开集上的限制&0，4.1.2\\
预概形态射在子预概形上的限制&I，4.1.7\\
预概形在开集上的限制&I，2.1.8\\
有理映射在开集上的限制&I，7.1.2\\
概形&I，5.4.1\\
仿射概形&I，1.7.1\\
局部概形；预概形在一点的局部概形&I，2.4.1\\
代数 $K$-概形&I，6.4.1\\
$K$ 上的有限概形；有限 $K$-概形&I，6.4.5\\
$S$-概形&I，5.4.1\\
仿射形式概形&I，10.1.2\\
$\frS$-形式概形&I，10.15.1\\
层的截面&0，3.1.6\\
$S$-预概形的 $S$-截面&I，2.5.5；5.3.11\\
$S$-预概形的有理 $S$-截面&I，7.1.2\\
$\OO_X$ 的单位截面&0，4.1.1\\
限制形式幂级数&0，7.5.1\\
预概形的和&I，3.6.1\\
由子模层生成的 $\OO_X$-子代数层&0，4.1.3\\
子预概形&I，4.1.3\\
与浸入相伴的子预概形&I，4.2.1\\
闭子预概形&I，4.1.3\\
由理想层定义的闭子预概形&I，4.1.2\\
闭形式子预概形&I，10.14.2\\
一点的特化&0，2.1.2\\
环的谱&I，1.1.1\\
可容许环的形式谱&I，10.1.2\\
模的支集&0，1.7.1\\
群层的支集&0，3.1.6\\
定义理想的基本系&I，10.3.6；10.5.1\\
预概形的进制归纳系&I，10.12.2\\
$J$-进制拓扑\textsuperscript{1}；$J$-预进制拓扑&0，7.1.9；7.2.3\\
谱拓扑&I，1.1.2\\
平凡的可逆 $\OO_X$-模层&I，2.4.8\\
有限型 $\OO_X$-模层&0，5.2.1\\
截面在一点的值&0，5.5.1\\
\end{longtable}
\noindent\textsuperscript{1}\textbf{原注：}注意，我们的术语与通常的定义略有不同：只有在所涉及的环或模\emph{分离且完备}时，我们才使用“进制”一词。
\endgroup


% END terminology_index.tex

\par

% BEGIN original_contents.tex
\chapter*{第一卷原著目录}
\addcontentsline{toc}{chapter}{第一卷原著目录}
\markboth{第一卷原著目录}{第一卷原著目录}
\src{225—228}
\noindent\textbf{译注。}本目录保留原书印刷页码，供与法文底本互查；本译稿的实际页码请见卷首目录。第一卷原著目录中的标题依原文译出，个别标题与正文略有差别的情形不另行统一。
\begingroup\small
\renewcommand{\arraystretch}{1.1}
\begin{longtable}{@{}p{0.84\linewidth}r@{}}
\toprule 标题&原书页码\\\midrule\endfirsthead
\toprule 标题（续）&原书页码\\\midrule\endhead
\bottomrule\endfoot
\textbf{导言}&5\\
\textbf{第 0 章\quad 预备知识}&11\\
\textbf{§1\quad 分式环}&11\\
1.0\quad 环与代数&11\\
1.1\quad 理想的根；环的幂零根与根&12\\
1.2\quad 分式模与分式环&13\\
1.3\quad 函子性质&14\\
1.4\quad 乘法子集的变更&15\\
1.5\quad 底环的变更&17\\
1.6\quad 将模 $M_f$ 等同于一个归纳极限&19\\
1.7\quad 模的支集&20\\
\textbf{§2\quad 不可约空间；诺特空间}&21\\
2.1\quad 不可约空间&21\\
2.2\quad 诺特空间&23\\
\textbf{§3\quad 层的补充知识}&23\\
3.1\quad 取值于一个范畴的层&23\\
3.2\quad 开集基上的预层&25\\
3.3\quad 层的粘合&28\\
3.4\quad 预层的直像&29\\
3.5\quad 预层的逆像&30\\
3.6\quad 常值层与局部常值层&33\\
3.7\quad 群预层或环预层的逆像&34\\
3.8\quad 伪离散空间层&35\\
\textbf{§4\quad 赋环空间}&35\\
4.1\quad 赋环空间、$\A$-模层、$\A$-代数层&35\\
4.2\quad $\A$-模层的直像&39\\
4.3\quad $\A$-模层的逆像&40\\
4.4\quad 直像与逆像的关系&42\\
\textbf{§5\quad 拟凝聚层与凝聚层}&44\\
5.1\quad 拟凝聚层&44\\
5.2\quad 有限型层&45\\
5.3\quad 凝聚层&47\\
5.4\quad 局部自由层&48\\
5.5\quad 局部环赋环空间上的层&53\\
\textbf{§6\quad 平坦性}&54\\
6.1\quad 平坦模&55\\
6.2\quad 底环的变更&55\\
6.3\quad 平坦性的局部化&56\\
6.4\quad 忠实平坦模&57\\
6.5\quad 标量限制&58\\
6.6\quad 忠实平坦环&58\\
6.7\quad 赋环空间的平坦态射&59\\
\textbf{§7\quad 进制环}&60\\
7.1\quad 可容许环&60\\
7.2\quad 进制环与射影极限&62\\
7.3\quad 诺特预进制环&66\\
7.4\quad 局部环上的拟有限模&68\\
7.5\quad 限制形式幂级数环&69\\
7.6\quad 完备分式环&72\\
7.7\quad 完备张量积&75\\
7.8\quad 同态模上的拓扑&77\\[0.3em]
\textbf{第 I 章\quad 概形的语言}&79\\
\textbf{§1\quad 仿射概形}&80\\
1.1\quad 环的素谱&80\\
1.2\quad 环素谱的函子性质&83\\
1.3\quad 与模相伴的层&84\\
1.4\quad 素谱上的拟凝聚层&90\\
1.5\quad 素谱上的凝聚层&92\\
1.6\quad 素谱上拟凝聚层的函子性质&93\\
1.7\quad 仿射概形态射的刻画&96\\
\textbf{§2\quad 预概形与预概形态射}&97\\
2.1\quad 预概形的定义&97\\
2.2\quad 预概形态射&98\\
2.3\quad 预概形的粘合&101\\
2.4\quad 局部概形&101\\
2.5\quad 预概形上的预概形&103\\
\textbf{§3\quad 预概形的积}&104\\
3.1\quad 预概形的和&104\\
3.2\quad 预概形的积&104\\
3.3\quad 积的形式性质；基预概形的变更&108\\
3.4\quad 预概形的预概形值点；几何点&111\\
3.5\quad 满射与单射&114\\
3.6\quad 纤维&117\\
3.7\quad 应用：预概形模 $J$ 的约化&118\\
\textbf{§4\quad 子预概形与浸入态射}&119\\
4.1\quad 子预概形&119\\
4.2\quad 浸入态射&122\\
4.3\quad 浸入的积&124\\
4.4\quad 预概形的逆像&125\\
4.5\quad 局部浸入与局部同构&126\\
\textbf{§5\quad 约化预概形；分离条件}&127\\
5.1\quad 约化预概形&127\\
5.2\quad 具有给定底空间的子预概形的存在性&131\\
5.3\quad 对角线；态射的图&132\\
5.4\quad 分离态射与分离预概形&135\\
5.5\quad 分离性判别准则&136\\
\textbf{§6\quad 有限性条件}&140\\
6.1\quad 诺特预概形与局部诺特预概形&140\\
6.2\quad 阿廷预概形&143\\
6.3\quad 有限型态射&144\\
6.4\quad 代数预概形&147\\
6.5\quad 态射的局部确定&150\\
6.6\quad 拟紧态射与局部有限型态射&152\\
\textbf{§7\quad 有理映射}&155\\
7.1\quad 有理映射与有理函数&155\\
7.2\quad 有理映射的定义域&158\\
7.3\quad 有理函数层&161\\
7.4\quad 挠层与无挠层&163\\
\textbf{§8\quad Chevalley 概形}&164\\
8.1\quad 相关的局部环&164\\
8.2\quad 整概形的局部环&165\\
8.3\quad Chevalley 概形&168\\
\textbf{§9\quad 拟凝聚层的补充知识}&169\\
9.1\quad 拟凝聚层的张量积&169\\
9.2\quad 拟凝聚层的直像&171\\
9.3\quad 拟凝聚层截面的延拓&172\\
9.4\quad 拟凝聚层的延拓&174\\
9.5\quad 预概形的闭像；子预概形的闭包&176\\
9.6\quad 拟凝聚代数层；结构层的变更&179\\
\textbf{§10\quad 形式概形}&180\\
10.1\quad 仿射形式概形&180\\
10.2\quad 仿射形式概形的态射&182\\
10.3\quad 仿射形式概形的定义理想&183\\
10.4\quad 形式预概形与形式预概形态射&185\\
10.5\quad 形式预概形的定义理想&186\\
10.6\quad 作为预概形归纳极限的形式预概形&188\\
10.7\quad 形式预概形的积&193\\
10.8\quad 预概形沿闭子集的形式完备化&194\\
10.9\quad 态射到完备化的延拓&198\\
10.10\quad 对仿射形式概形上凝聚层的应用&201\\
10.11\quad 形式预概形上的凝聚层&204\\
10.12\quad 形式预概形的进制态射&206\\
10.13\quad 有限型态射&207\\
10.14\quad 形式预概形的闭子预概形&209\\
10.15\quad 分离形式预概形&212\\
\textbf{参考文献}&215\\
\textbf{记号索引}&217\\
\textbf{术语索引}&219\\
\end{longtable}
\endgroup


% END original_contents.tex

\par

% BEGIN colophon.tex
\chapter*{第一卷书末出版记载}
\addcontentsline{toc}{chapter}{第一卷书末出版记载}
\markboth{第一卷书末出版记载}{第一卷书末出版记载}
\src{228}
收稿日期：1959 年 10 月 17 日。

\medskip
1960 年。法国大学出版社印刷厂（Imprimerie des Presses Universitaires de France）。法国旺多姆（Vendôme）。

\medskip
版号：25937。法国印制。印刷编号：16186。

\bigskip
\noindent\textbf{译注。}以上为法文原著的出版记载，不是本中文译稿的出版或审定信息。

\bigskip
\noindent\textbf{原扉页与版权页的出版信息。}
原扉页列明：高等科学研究所；《代数几何基础》，A. Grothendieck 著，
与 J. Dieudonné 合作撰写；《概形的语言》；1960 年；《数学出版物》第 4 号。
地址为法国巴黎第十六区 Bugeaud 环岛 5 号。

\medskip
原版权页载：法定送存；第一版，1960 年第三季度；在所有国家保留一切权利；
\copyright\ 1960，高等科学研究所（Institut des Hautes Études Scientifiques）。


% END colophon.tex

\par

% BEGIN copyright.tex
\chapter{第一卷底本的档案与版权说明}
\markboth{档案与版权说明}{档案与版权说明}
以下中文为上传文件首页所附档案说明的译文；随后照录其法文版权声明。

\medskip
\noindent\textbf{文献信息}

ALEXANDER GROTHENDIECK\par
《代数几何基础》（与 Jean Dieudonné 合作撰写）：I．概形的语言。\par
\textit{Publications mathématiques de l'I.H.É.S.}，第 4 卷（1960），5–228 页。\par
档案条目标识：\texttt{PMIHES\_1960\_\_4\_\_5\_0}。

\medskip
\noindent\textbf{档案说明译文}

\copyright\ \textit{Publications mathématiques de l'I.H.É.S.}，1960，保留所有权利。
访问本刊档案，即表示同意其一般使用条件。
原档案声明：任何商业使用或系统性印刷均构成刑事侵权；本文件的任何复制或打印都必须包含这一版权声明。
本文由“古旧数学文献数字化”（Numérisation de documents anciens mathématiques）计划数字化。

\medskip
\noindent\textbf{法文版权声明照录}

\begingroup\small
\noindent
\copyright\ Publications mathématiques de l'I.H.É.S., 1960, tous droits réservés.

L'accès aux archives de la revue « Publications mathématiques de l'I.H.É.S. »
(\url{http://www.ihes.fr/IHES/Publications/Publications.html}), implique l'accord avec les conditions générales
 d'utilisation (\url{http://www.numdam.org/legal.php}).
Toute utilisation commerciale ou impression systématique est constitutive d'une infraction pénale.
Toute copie ou impression de ce fichier doit contenir la présente mention de copyright.

Article numérisé dans le cadre du programme\par
Numérisation de documents anciens mathématiques\par
\url{http://www.numdam.org/}
\endgroup

\medskip
\noindent\textbf{版本说明}

本书为根据用户所上传原文制作的全卷中文译稿，并非 IHÉS、NUMDAM 或原权利人发布的中文版本。
原版中的历史陈述、书目年代与引文编号均按源文件保留。


% END copyright.tex

